% Осень26!B18; нумерация внутри столбца без пустых строк.
% Источник: лекция 3, §4 (строка 115)
\definitionticket{15}{Полная мера}
\label{definition:15}

\subsection*{§ «Продолжение меры»: полная и σ-конечная меры}

Дальше --- движение «в сторону приложений: лично меня интересует мера Лебега, вас, наверное, ещё и вероятностные приложения».

\begin{definition}

\emph{полная мера.} Пусть \(\mathcal{A}\) --- σ-алгебра, \(\mu:\mathcal{A}\to\overline{\mathbb{R}}\) --- мера. \(\mu\) называется \textbf{полной}, если \[\forall A\in\mathcal{A}\ \text{с}\ \mu A=0\quad\forall B\subset A:\qquad B\in\mathcal{A}\] (а тогда автоматически и \(\mu B=0\) --- по монотонности меры).

\end{definition}

\textbf{Комментарии лектора.}

\begin{itemize}
\tightlist
\item
  «Фактически я рассказываю не про меру, а про σ-алгебру»: полнота --- свойство σ-алгебры вместе с каждым множеством меры ноль содержать \textbf{все} его подмножества. Множество меры ноль потенциально континуально, и тогда σ-алгебра должна содержать \(2^{\mathfrak{c}}\) элементов --- «количественно большая».
\item
  Здравый смысл: «если уж что-то померили и оно имеет меру ноль, то как можно думать, что его меньшая часть неизмерима? Должно быть измеримо». Но это требует проверки конструкции.
\end{itemize}
